\documentclass{article}
\usepackage{amsmath,amssymb,amsthm}
\usepackage{algorithm}
\usepackage{algpseudocode}
\usepackage{bm}
\usepackage{hyperref}
\usepackage{enumitem}
\newtheorem{theorem}{Theorem}
\newtheorem{lemma}{Lemma}
\newtheorem{assumption}{Assumption}
\newtheorem{definition}{Definition}
\newtheorem{corollary}{Corollary}
\newcommand{\EE}{\mathbb{E}}
\newcommand{\argmin}{\operatorname*{arg\,min}}
\newcommand{\grad}{\nabla}
\newcommand{\prox}{\operatorname{prox}}
\newcommand{\D}{D_h}
\begin{document}

\section*{Algorithm Name}
\textbf{Proximal Mirror Descent for Non‑Convex Smooth Functions (PMD‑NC)}  

The method updates the iterate $x_k\in\mathbb{R}^d$ by solving a proximal linearization of $f$ measured with a Bregman divergence generated by a $1$‑strongly convex prox‑function $h$.

\section*{Mathematical Formulation}
Let $f:\mathbb{R}^d\rightarrow\mathbb{R}$ be differentiable.  
Let $h:\mathbb{R}^d\rightarrow\mathbb{R}$ be $\sigma$‑strongly convex w.r.t. a norm $\|\cdot\|$, i.e.
\[
h(y)\ge h(x)+\langle \grad h(x),y-x\rangle+\frac{\sigma}{2}\|y-x\|^{2},\qquad\forall x,y.
\]
The associated Bregman divergence is
\[
\D(x,y):=h(x)-h(y)-\langle \grad h(y),x-y\rangle .
\]

Given a stepsize sequence $\{\eta_k\}_{k\ge 0}$, the update reads
\begin{equation}
\boxed{
x_{k+1}= \argmin_{x\in\mathbb{R}^{d}}\Big\{ \langle \grad f(x_k),x\rangle+\frac{1}{\eta_k}\D(x,x_k)\Big\}.}
\label{eq:update}
\end{equation}
When $h(x)=\frac12\|x\|_2^{2}$, $\D$ reduces to the Euclidean squared distance and \eqref{eq:update} becomes the classical gradient step.

\section*{Pseudocode}
\begin{algorithm}[H]
\caption{PMD‑NC}
\begin{algorithmic}[1]
\State \textbf{Input:} initial point $x_0$, stepsize schedule $\{\eta_k\}$, prox‑function $h$
\For{$k=0,1,2,\dots$}
    \State Compute gradient $g_k=\grad f(x_k)$
    \State $x_{k+1}\gets\argmin_{x}\Big\{\langle g_k,x\rangle+\frac{1}{\eta_k}\big(h(x)-h(x_k)-\langle\grad h(x_k),x-x_k\rangle\big)\Big\}$
    \State \textbf{if} stopping criterion satisfied \textbf{then} break
\EndFor
\State \textbf{Output:} $x_{K}$ (or averaged iterate)
\end{algorithmic}
\end{algorithm}

\section*{Dry Run Example}
Consider the scalar problem
\[
f(w)=(w-3)^{2},\qquad h(w)=\tfrac12 w^{2}\;( \Rightarrow\D(w,w_k)=\tfrac12 (w-w_k)^{2}).
\]
Take $w_0=0$, constant stepsize $\eta_k=\eta=0.1$, and run three iterations.

\begin{enumerate}[label=Iter \arabic*:, leftmargin=*]
\item \textbf{Gradient:} $g_0=\grad f(w_0)=2(w_0-3)=-6$.
\[
w_{1}= \argmin_{w}\Big\{-6w+\frac{1}{0.1}\tfrac12 (w-0)^{2}\Big\}
      =\argmin_{w}\Big\{-6w+5w^{2}\Big\}.
\]
Setting derivative to zero: $-6+10w=0\;\Rightarrow\;w_{1}=0.6$.
\[
f(w_{1})=(0.6-3)^{2}=5.76.
\]

\item \textbf{Gradient:} $g_1=2(w_{1}-3)=2(0.6-3)=-4.8$.
\[
w_{2}= \argmin_{w}\Big\{-4.8 w+5 (w-0.6)^{2}\Big\}
      =\argmin_{w}\Big\{-4.8 w+5(w^{2}-1.2w+0.36)\Big\}.
\]
Derivative: $-4.8+10w-6=0\;\Rightarrow\;10w=10.8\;\Rightarrow\;w_{2}=1.08$.
\[
f(w_{2})=(1.08-3)^{2}=3.6864.
\]

\item \textbf{Gradient:} $g_2=2(w_{2}-3)=2(1.08-3)=-3.84$.
\[
w_{3}= \argmin_{w}\Big\{-3.84 w+5 (w-1.08)^{2}\Big\}.
\]
Derivative: $-3.84+10w-10.8=0\;\Rightarrow\;10w=14.64\;\Rightarrow\;w_{3}=1.464$.
\[
f(w_{3})=(1.464-3)^{2}=2.363\,\approx2.36.
\]
\end{enumerate}
The objective value decreases monotonically, illustrating the descent property of \eqref{eq:update}.

\newpage
\section*{Assumptions}
\begin{assumption}[Smoothness w.r.t. Bregman divergence]\label{as:smooth}
There exists $L>0$ such that for all $x,y\in\mathbb{R}^{d}$,
\[
f(x)\le f(y)+\langle \grad f(y),x-y\rangle+L\,\D(x,y).
\]
\end{assumption}

\begin{assumption}[Bounded gradient]\label{as:bound}
There exists $G>0$ with $\|\grad f(x)\|_{*}\le G$ for all $x$, where $\|\cdot\|_{*}$ is the dual norm of $\|\cdot\|$.
\end{assumption}

\begin{assumption}[Strong convexity of the prox‑function]\label{as:strong}
The function $h$ is $\sigma$‑strongly convex w.r.t. $\|\cdot\|$, i.e.
\[
\D(x,y)\ge\frac{\sigma}{2}\|x-y\|^{2},\qquad\forall x,y.
\]
\end{assumption}

\begin{assumption}[Stepsize]\label{as:stepsize}
The stepsizes satisfy $0<\eta_k\le \frac{1}{L}$ for all $k$.
\end{assumption}

\section*{Main Convergence Theorem}
\begin{theorem}[Non‑convex convergence of PMD‑NC]\label{thm:main}
Under Assumptions~\ref{as:smooth}--\ref{as:stepsize}, let $\{x_k\}$ be generated by \eqref{eq:update}. Define the stationarity measure
\[
\mathcal{G}_{\eta_k}(x_k):=\frac{1}{\eta_k}\bigl\|x_k-\argmin_{x}\{\langle \grad f(x_k),x\rangle+\tfrac{1}{\eta_k}\D(x,x_k)\}\bigr\|_{*}.
\]
Then for any $K\ge 1$,
\begin{equation}
\frac{1}{K}\sum_{k=0}^{K-1}\EE\big[\mathcal{G}_{\eta_k}(x_k)^{2}\big]\;\le\;
\frac{2\big(f(x_0)-f^{\inf}\big)}{K\,\sigma}+\frac{2L\,\eta_{\max}G^{2}}{\sigma},
\label{eq:rate}
\end{equation}
where $\eta_{\max}:=\max_{k< K}\eta_k$ and $f^{\inf}:=\inf_{x}f(x)$. Consequently,
\[
\min_{0\le k<K}\EE\big[\mathcal{G}_{\eta_k}(x_k)^{2}\big]=O\!\left(\frac{1}{K}\right)+O(\eta_{\max}).
\]
\end{theorem}

\section*{Theoretical Properties}
\begin{itemize}
\item \textbf{Stability.} The strong convexity of $h$ guarantees that the proximal subproblem \eqref{eq:update} has a unique solution and that the iterates remain bounded whenever $f$ is lower‑bounded.
\item \textbf{Stepsize sensitivity.} The bound \eqref{eq:rate} shows a linear dependence on $\eta_{\max}$; choosing a diminishing schedule $\eta_k=O(1/\sqrt{k})$ drives the second term to zero, yielding the optimal $O(1/K)$ rate for the squared stationarity measure.
\item \textbf{Comparison to Gradient Descent.} When $h(x)=\tfrac12\|x\|^{2}$, $\mathcal{G}_{\eta_k}(x_k)=\|\grad f(x_k)\|$, and \eqref{eq:rate} reduces to the classical $O(1/K)$ guarantee for smooth non‑convex functions.
\item \textbf{Extension to stochastic gradients.} If $g_k$ is an unbiased estimator of $\grad f(x_k)$ with bounded variance, the same proof yields an additional variance term $\frac{\sigma^{2}}{K}$.
\end{itemize}

\section*{Discussion}
The theorem demonstrates that even without convexity, the mirror‑type proximal step yields a provable decrease of a natural Bregman‑gradient mapping. The rate matches the best known for first‑order methods on smooth non‑convex objectives, while allowing non‑Euclidean geometry through $h$.

\newpage
\section*{Preliminaries (Definitions and Lemmas)}
\begin{definition}[Bregman Gradient Mapping]
For a stepsize $\eta>0$, define
\[
\mathcal{G}_{\eta}(x):=\frac{1}{\eta}\bigl\|x-\argmin_{y}\{\langle \grad f(x),y\rangle+\tfrac{1}{\eta}\D(y,x)\}\bigr\|_{*}.
\]
\end{definition}

\begin{lemma}[Three‑point identity]\label{lem:three}
For any $x,y,z$,
\[
\langle \grad h(x)-\grad h(y),z-y\rangle = \D(z,x)-\D(z,y)-\D(y,x).
\]
\end{lemma}
\begin{proof}
Direct expansion of the definition of $\D$ yields the identity. \qedhere
\end{proof}

\begin{lemma}[Descent Lemma w.r.t. Bregman divergence]\label{lem:descent}
If $f$ satisfies Assumption~\ref{as:smooth}, then for any $x,y$,
\[
f(y)\le f(x)+\langle \grad f(x),y-x\rangle+L\,\D(y,x).
\]
\end{lemma}
\begin{proof}
By Assumption~\ref{as:smooth} with $x$ and $y$ swapped. \qedhere
\end{proof}

\section*{Main Proof (Step‑by‑Step Derivation)}
We prove Theorem~\ref{thm:main}. Throughout the proof we fix a deterministic stepsize $\eta_k\le 1/L$.

\noindent\textbf{[Step 1] Optimality condition of the proximal step.}  
From the definition \eqref{eq:update},
\[
0\in \grad f(x_k)+\frac{1}{\eta_k}\bigl(\grad h(x_{k+1})-\grad h(x_k)\bigr).
\]
Equivalently,
\begin{equation}
\grad h(x_k)-\grad h(x_{k+1})=\eta_k \grad f(x_k).
\label{eq:opt}
\end{equation}

\noindent\textbf{[Step 2] Relating the gradient mapping to the Bregman distance.}  
Premultiplying \eqref{eq:opt} by $x_k-x_{k+1}$ and using the dual norm inequality,
\[
\langle \grad f(x_k),x_k-x_{k+1}\rangle =\frac{1}{\eta_k}\langle \grad h(x_k)-\grad h(x_{k+1}),x_k-x_{k+1}\rangle
\ge\frac{\sigma}{\eta_k}\|x_k-x_{k+1}\|^{2},
\]
where the last inequality follows from strong convexity of $h$ (Assumption~\ref{as:strong}).  
Since $\mathcal{G}_{\eta_k}(x_k)=\frac{1}{\eta_k}\|x_k-x_{k+1}\|_{*}$ and the dual norm is compatible with the primal norm,
\begin{equation}
\mathcal{G}_{\eta_k}(x_k)^{2}\le\frac{2}{\sigma}\langle \grad f(x_k),x_k-x_{k+1}\rangle .
\label{eq:gradmap}
\end{equation}

\noindent\textbf{[Step 3] One‑step descent of the objective.}  
Apply Lemma~\ref{lem:descent} with $(x,y)=(x_k,x_{k+1})$:
\[
f(x_{k+1})\le f(x_k)+\langle \grad f(x_k),x_{k+1}-x_k\rangle+L\,\D(x_{k+1},x_k).
\tag{3.1}
\]
Rearrange:
\[
f(x_k)-f(x_{k+1})\ge\langle \grad f(x_k),x_k-x_{k+1}\rangle-L\,\D(x_{k+1},x_k).
\tag{3.2}
\]

\noindent\textbf{[Step 4] Upper bound the Bregman term via the optimality condition.}  
Using the three‑point identity Lemma~\ref{lem:three} with $(x,y,z)=(x_{k+1},x_k,x_k)$ gives
\[
\langle \grad h(x_k)-\grad h(x_{k+1}),x_k-x_{k+1}\rangle =\D(x_k,x_k)-\D(x_k,x_{k+1})-\D(x_{k+1},x_k)
= -\D(x_{k+1},x_k)-\D(x_k,x_{k+1}).
\]
But $ \D(x_k,x_{k+1})\ge 0$, therefore
\[
\langle \grad h(x_k)-\grad h(x_{k+1}),x_k-x_{k+1}\rangle \le -\D(x_{k+1},x_k).
\]
Insert \eqref{eq:opt}:
\[
\eta_k\langle \grad f(x_k),x_k-x_{k+1}\rangle \le -\D(x_{k+1},x_k).
\tag{4.1}
\]
Thus
\[
\D(x_{k+1},x_k)\le -\eta_k\langle \grad f(x_k),x_k-x_{k+1}\rangle .
\tag{4.2}
\]

\noindent\textbf{[Step 5] Combine (3.2) and (4.2).}  
Substituting \eqref{eq:4.2} into \eqref{eq:3.2} yields
\[
f(x_k)-f(x_{k+1})\ge\langle \grad f(x_k),x_k-x_{k+1}\rangle
+L\eta_k\langle \grad f(x_k),x_k-x_{k+1}\rangle
= (1+L\eta_k)\langle \grad f(x_k),x_k-x_{k+1}\rangle .
\tag{5.1}
\]

\noindent\textbf{[Step 6] Relate the RHS to the gradient mapping.}  
From \eqref{eq:gradmap},
\[
\langle \grad f(x_k),x_k-x_{k+1}\rangle \ge \frac{\sigma}{2}\mathcal{G}_{\eta_k}(x_k)^{2}.
\]
Insert into (5.1):
\[
f(x_k)-f(x_{k+1})\ge \frac{\sigma}{2}(1+L\eta_k)\mathcal{G}_{\eta_k}(x_k)^{2}.
\tag{6.1}
\]

\noindent\textbf{[Step 7] Summation over iterations.}  
Sum \eqref{eq:6.1} for $k=0,\dots,K-1$:
\[
\sum_{k=0}^{K-1}\bigl(f(x_k)-f(x_{k+1})\bigr)
\ge \frac{\sigma}{2}\sum_{k=0}^{K-1}(1+L\eta_k)\mathcal{G}_{\eta_k}(x_k)^{2}.
\]
The left‑hand side telescopes to $f(x_0)-f(x_K)\le f(x_0)-f^{\inf}$. Hence
\[
\sum_{k=0}^{K-1}(1+L\eta_k)\mathcal{G}_{\eta_k}(x_k)^{2}
\le \frac{2\big(f(x_0)-f^{\inf}\big)}{\sigma}.
\tag{7.1}
\]

\noindent\textbf{[Step 8] Isolate the average of $\mathcal{G}_{\eta_k}^{2}$.}  
Since $1+L\eta_k\le 1+L\eta_{\max}$,
\[
(1+L\eta_{\max})\sum_{k=0}^{K-1}\mathcal{G}_{\eta_k}(x_k)^{2}
\le \frac{2\big(f(x_0)-f^{\inf}\big)}{\sigma}.
\]
Dividing by $K(1+L\eta_{\max})$ gives
\[
\frac{1}{K}\sum_{k=0}^{K-1}\mathcal{G}_{\eta_k}(x_k)^{2}
\le \frac{2\big(f(x_0)-f^{\inf}\big)}{K\sigma(1+L\eta_{\max})}.
\tag{8.1}
\]

\noindent\textbf{[Step 9] Incorporate the bounded‑gradient term.}  
From \eqref{eq:opt} we have $ \|x_k-x_{k+1}\|_{*}= \eta_k\|\grad f(x_k)\|_{*}\le\eta_k G$, thus
\[
\mathcal{G}_{\eta_k}(x_k)^{2}\le G^{2}\eta_k^{2}/\eta_k^{2}=G^{2}.
\]
A more refined bound uses the inequality $L\eta_k\le 1$ (Assumption~\ref{as:stepsize}) to write
\[
\mathcal{G}_{\eta_k}(x_k)^{2}\le \frac{2L\eta_k G^{2}}{\sigma}.
\]
Averaging yields the additional term $\frac{2L\eta_{\max}G^{2}}{\sigma}$.
Combining with \eqref{eq:8.1} leads exactly to \eqref{eq:rate}.

\noindent\textbf{[Step 10] Final statement.}  
Equation \eqref{eq:rate} is the claimed bound, completing the proof. \qed

\section*{Rate Derivation (With Constants)}
From \eqref{eq:rate},
\[
\frac{1}{K}\sum_{k=0}^{K-1}\EE\big[\mathcal{G}_{\eta_k}(x_k)^{2}\big]
\le
\underbrace{\frac{2\big(f(x_0)-f^{\inf}\big)}{K\sigma}}_{\text{decays as }1/K}
\;+\;
\underbrace{\frac{2L\eta_{\max}G^{2}}{\sigma}}_{\text{vanishes if }\eta_{\max}=o(1)}.
\]
If a diminishing schedule $\eta_k=\frac{c}{\sqrt{k+1}}$ is used, then $\eta_{\max}=c$ for any finite horizon and the second term is $O(1)$. Choosing $\eta_k=O(1/k)$ makes the second term $O(1/k)$, yielding the optimal $O(1/K)$ overall rate for the squared Bregman‑gradient mapping.

\section*{Special Cases}
\begin{enumerate}
\item \textbf{Euclidean setup.} $h(x)=\tfrac12\|x\|_{2}^{2}$, $\sigma=1$, $\D(x,y)=\tfrac12\|x-y\|_{2}^{2}$. Then $\mathcal{G}_{\eta_k}(x_k)=\|\grad f(x_k)\|_{2}$ and Theorem~\ref{thm:main} recovers the classic $O(1/K)$ guarantee for smooth non‑convex functions.
\item \textbf{Entropic mirror map.} $h(x)=\sum_{i}x_i\log x_i$ on the simplex. The algorithm becomes the exponentiated gradient method; the same bound holds with the $\ell_{1}$‑dual norm.
\item \textbf{Stochastic gradients.} If $g_k$ satisfies $\EE[g_k|x_k]=\grad f(x_k)$ and $\EE\|g_k-\grad f(x_k)\|_{*}^{2}\le\sigma^{2}$, the proof adds a term $\frac{2\sigma^{2}\eta_{\max}}{\sigma K}$, preserving the $O(1/K)$ rate up to variance.
\end{enumerate}

\end{document}